In the work, of which the present volume is an instalment, my endeavour
has been to lay before the reader a connected exposition of the theory
of sound, which should include the more important of the advances made
in modern times by Mathematicians and Physicists. The importance of the
object which I have had in view will not, I think, be disputed by those
competent to judge. At the present time many of the most valuable
contributions to science are to be found only in scattered periodicals
and transactions of societies, published in various parts of the world
and in several languages, and are often practically inaccessible to
those who do not happen to live in the neighbourhood of large public
libraries. In such a state of things the mechanical impediments to study
entail an amount of unremunerative labour and consequent hindrance to
the advancement of science which it would be difficult to overestimate.

Since the well-known Article on Sound in the Encyclopedia Metropolitana,
by Sir John Herschel (1845), no complete work has been published in
which the subject is treated mathematically. By the premature death of
Prof. Donkin the scientific world was deprived of one whose mathematical
attainments in combination with a practical knowledge of music qualified
him in a special manner to write on Sound. The first part of his
Acoustics (1870), though little more than a fragment, is sufficient to
shew that my labours would have been unnecessary had Prof. Donkin lived
to complete his work.

In the choice of topics to be dealt with in a work on Sound, I have for
the most part followed the example of my predecessors. To a great extent
the theory of Sound, as commonly understood, covers the same ground as
the theory of Vibrations in general ; but, unless some limitation were
admitted, the consideration of such subjects as the Tides, not to speak
of Optics, would have to be included. As a general rule we shall confine
ourselves to those classes of vibrations for which our ears afford a
ready made and wonderfully sensitive instrument of investigation.
Without ears we should hardly care much more about vibrations than
without eyes we should care about light.

The present volume includes chapters on the vibrations of systems in
general, in which, I hope, will be recognised some novelty of treatment
and results, followed by a more detailed consideration of special
systems, such as stretched strings, bars, membranes, and plates. The
second volume, of which a considerable portion is already written, will
commence with aerial vibrations.

My best thanks are due to Mr H. M. Taylor of Trinity College, Cambridge,
who has been good enough to read the proofs. By his kind assistance
several errors and obscurities have been eliminated, and the volume
generally has been rendered less imperfect than it would otherwise have
been.

Any corrections, or suggestions for improvements, with which my readers
may favour me will be highly appreciated.

Terling Place, Witham, April, 1877,


IN this \textbf{second} edition all corrections of importance are noted, and new
matter appears either as fresh sections, e.g. § 32 a, or enclosed in
square brackets {[} {]}. Two new chapters X a, X B are interpolated,
devoted to Curved Plates or Shells, and to Electrical Vibrations Much of
the additional matter relates to the more difficult parts of the subject
and will be passed over by the reader on a first perusal.

In the mathematical investigations I have usually employed such methods
as present themselves naturally to a physicist. The pure mathematician
will complain, and (it must be confessed) sometimes with justice, of
deficient rigour. But to the question there are two sides. For, however
important it may be to maintain a uniformly high standard in pure
mathematics, the physicist may occasionally do well to rest content with
arguments which are fairly satisfactory and conclusive from his point
of view. To his mind, exercised in a different order of ideas, the more
severe procedure of the pure mathematician may appear not more but less
demonstrative. And further, in many cases of difficulty to insist
upon the highest standard would mean the exclusion of the subject
altogether in view of the space that would be required.

In the first edition much stress was laid upon the establishment of
general theorems by means of Lagrange's method, and I am more than ever
impressed with the advantages of this procedure. It not infrequently
happens that a theorem can.be thus demonstrated in all its generality
with less mathematical apparatus than is required for dealing with
particular cases by special methods.

During the revision of the proof-sheets I have again had the very great
advantage of the cooperation of Mr H. M. Taylor, until he was
unfortunately compelled to desist. To him and to several other friends
my thanks are due for valuable suggestions.

J W S, 1894.

EDITORIAL NOTE.

In this reissue, a few pencilled corrections and references in the
Author's own copy have been made use of. Otherwise no change has been made.
